\section{2030 Reapportionment Projections}
\label{sec:projections}

\subsection{Projection Scenario}

Based on the projection scenario used in this paper~\citep{census2030proj}
and \hh{} apportionment calculations~\citep{huntington1928,balinski1982},
we analyze the following possible seat changes for 2030:

\begin{table}[h]
\centering
\caption{Illustrative 2030 congressional seat-change scenario. Seat counts
         shown as 2020 $\to$ 2030 scenario value. ``Tree change'' indicates
         the qualitative change in the \ar{} split tree.}
\label{tab:projections}
\small
\begin{tabular}{lcccc}
\toprule
State & 2020 seats & 2030 (proj.) & Change & Tree change \\
\midrule
Texas & 38 & 41 & $+3$ & Dramatic (19-way $\to$ binary fallback) \\
California & 52 & 50 & $-2$ & Structural \\
Florida & 28 & 29 & $+1$ & Dramatic (7-way $\to$ binary fallback) \\
New York & 26 & 25 & $-1$ & Structural \\
Arizona & 9 & 10 & $+1$ & Minor \\
Georgia & 14 & 15 & $+1$ & Minor \\
North Carolina & 14 & 15 & $+1$ & Minor \\
Illinois & 17 & 16 & $-1$ & Minor \\
Michigan & 13 & 12 & $-1$ & Moderate \\
Pennsylvania & 17 & 16 & $-1$ & Minor \\
\bottomrule
\end{tabular}
\end{table}

\subsection{The Role of Prime Factorization}

The implemented \ar{} split prescription for $k$ districts is constructed as
follows:
\begin{enumerate}
\item Find the prime factorisation $k = p_1^{a_1} \times \cdots \times p_r^{a_r}$.
\item If $k \leq 3$, perform a direct $k$-way split.
\item If $k$ is composite, split by the largest prime factor $p$, producing
      $p$ subregions with $k/p$ districts each.
\item If $k > 3$ is prime, split into two subregions with
      $\lfloor k/2 \rfloor$ and $\lceil k/2 \rceil$ districts, then recurse.
\end{enumerate}

A change from composite $k$ to prime $k'$ fundamentally changes the top
prescription: the composite case starts with a largest-prime multiway split,
while the prime case starts with an uneven binary split.
This is the Texas and Florida case.  The old draft described large primes as
flat $k$-way splits; the current implementation deliberately avoids that by
using the binary fallback.

\subsection{Show the Tree, Not Just the Factorisation}

The difference is easiest to see as a small tree sketch:
\[
\begin{array}{c@{\qquad}c}
\text{Texas 2020: }38=2\times 19 &
\boxed{38}\rightarrow 19\times\boxed{2}\\[0.4em]
\text{Texas scenario: }41\text{ prime} &
\boxed{41}\rightarrow \boxed{20}+\boxed{21}
       \rightarrow (5\times\boxed{4}) + (7\times\boxed{3})
\end{array}
\]
The 2020 tree begins by carving 19 equal two-seat containers.  The scenario
tree begins by carving one 20-seat side and one 21-seat side, then decomposes
those subproblems.  The visible reason for disruption is therefore not merely
``prime versus composite''; it is a different first split of the state.

\subsection{Projection Uncertainty}

The 2030 seat counts are scenario values, not final values.
Seat changes can move at the margin, especially for fast-growing or
slow-declining states.
The structural conclusion should therefore be read conditionally:
if Florida is 29, it uses the large-prime binary fallback; if Florida is 30,
it becomes a composite $2\times3\times5$ case with a largest-prime first split.
The projection package should include a $\pm1$ sensitivity table before the
paper makes strong state-specific claims.
